\documentclass[10pt,a4paper]{amsart}
\usepackage[margin=19mm]{geometry}
\usepackage{amsmath,amssymb,mathtools}
\usepackage[scr=rsfs,cal=euler]{mathalfa}
\usepackage{xfrac}
\usepackage[unicode,hidelinks,bookmarks=false]{hyperref}
\newcommand{\ie}{i.e.\ }
\newcommand{\cf}{cf.\ }
\newcommand{\resp}{respectively\ }
\newcommand{\Subsection}{\subsection}
\providecommand{\nobreakdash}{\nobreak\hbox{-}}
\setlength{\emergencystretch}{2em}
\multlinegap0pt
\usepackage{amssymb}
\usepackage{mathtools}
\usepackage{float}
\usepackage{algorithm}\usepackage{algpseudocode}
\usepackage{tikz}
\newtheorem{theorem}{Theorem}
\newcommand{\setc}{\mathcal{C}}
\newcommand{\sete}{\mathcal{E}}
\newcommand{\setv}{\mathcal{V}}
\title{Construction and Decoding of Quantum Margulis Codes}
\author{Michele Pacenti, Dimitris Chytas, Bane Vasic}
\date{Source: arXiv:2503.03936v3 (CC BY 4.0)}
\begin{document}
\maketitle

\noindent\emph{Source and licence.} Construction and Decoding of Quantum Margulis Codes. Version arXiv:2503.03936v3 (CC BY 4.0), distributed by arXiv under the Creative Commons Attribution 4.0 International licence, http://creativecommons.org/licenses/by/4.0/, as recorded in the arXiv licence metadata for this exact version and shown on the abstract page https://arxiv.org/abs/2503.03936v3. Full licence notice: \texttt{licenses/arxiv-2503.03936-CC-BY-4.0.txt}.\par\smallskip

\noindent\textit{Editorial scope.} Counted unit: the article's theorem th:invariance (source0.tex lines 266--272). The article's complete original proof of it is delivered with it (source0.tex lines 276--278, one \texttt{proof} environment of 3 lines, which the article places away from the statement). No mathematics is written, repaired or re-derived: the statement and the proof are the article's own lines, in the article's own notation, and no retained line is substituted or re-flowed.  No further source-local context is needed: every cross-reference of every retained line resolves inside the two delivered spans.  Nothing was omitted from any retained span, so the package carries no omission record.  No retained line contains a citation, so the wrapper carries no bibliography and no bibliography entry.  No retained line contains a figure environment, an image-inclusion command or drawing code, so no image is delivered and the package declares no asset dependency.  Editorial formatting only, in addition: an AMS \texttt{amsart} preamble that restates the article's own declarations and macros used by the retained lines, namely the environments \texttt{theorem} on separate counters, together with the standard packages those lines need.  The retained environments print in this document as Theorem 1 in order of appearance, whereas the article prints the same items with the numbering of its own full document.  The complete native source of the article as distributed in the arXiv e-print for this exact version is bundled unchanged as \texttt{sources/174295/source0.tex}.
\begin{theorem}
\label{th:invariance}
    Let $\mathcal{T} = (\setv \cup \setc,\sete)$ be the $X$ (or $Z$) Tanner graph of a 2BGA code obtained from a group $G$ and two sets of generators $A,B$. Then, if $A$ ($B$) is a normal subgroup of $G$, $G \subseteq \mathrm{Aut}(\mathcal{T})$, assuming it to act on the right (left).
    \begin{proof}
        Consider a check node labeled as $g \in G$ and the two edges $(g,ga)$ and $(g,bg)$. If $\mathcal{T}$ is $G$-invariant, assuming $G$ to act on the right, for any $h\in G$ we have that the edges $(gh,gah)$ and $(gh,bgh)$ must exist. Clearly, since the check $gh$ must exist, the edge $(gh,bgh)$ is obtained by multiplying $gh$ by $b$ on the left, therefore such edge exists. On the other hand, if $(gh,gah)$ exists it means there exists $x\in A$ such that $ghx = gah$, meaning that $x=h^{-1}ah$. If $A$ is a normal subgroup of $G$, it is invariant under conjugation, \textit{i.e.}, for any $h \in G$ we have $h^{-1}Ah=A$, thus $x\in A$ and the edge exists. Assuming the action of $G$ to be from the right, the proof is analogous and it requires $B$ to be normal.
    \end{proof}
\end{theorem}

    \begin{proof}
    If $G$ is Abelian, all its subgroups are normal, therefore the hypothesis for Theorem \ref{th:invariance} is always satisfied.
    \end{proof}

\end{document}
